% Ordinary mathematical proof for the candidate PA194F.
% Requires amsmath, amssymb, amsthm, booktabs, and theorem/lemma environments.
\section{The new source transformation and its PA proof}
\label{sec:pa194f-proof}

The new machine PA194F retains the symmetric cached controller of the
supplied PA194CS machine and changes the last source instruction of the
axiom-emission command. Its only native transition change is
\[
                  (2066,1):0L302\quad\longmapsto\quad1R21.
\]
The working-state count remains 194. This is a change to the represented
instruction and its physical decoder, not a renaming of controls: an
unconditional source jump becomes a decrement with two different source
successors. The extra branch is unreachable from every intended source
start. The proof below establishes that assertion by an invariant involving
only four program locations.


The general zero-closing idea is inherited. The supplied 196-state
research already replaced the last AX break by a decrement of a register
known to be zero, with separate variants for registers 3, 4, 5 and 6;
see \texttt{TRIT\_CONTEXTS.md} in the earlier comparison-and-new196
packet~\cite{closing}. The new result here is its complete verified composition with
the 194-state symmetric cached controller and its particular decoder,
together with the four-site proof below. Novelty is claimed for the
resulting literal control graph within the pinned corpus, not for the
general principle of replacing a jump by a known-zero test.

\subsection{The exact four-site invariant}

Let $P$ be the 3,646-site source program of PA194CS and let $P_F$ be the
source of PA194F. The programs have the same instruction at every site
except 3553. A decrement
$\operatorname{dec}r[z,p]$ leaves a zero register unchanged and goes to $z$;
on a positive register it subtracts one and goes to $p$. Put $X=r_6$.
The four sites in
\[
                         Z=\{3550,3551,3552,3553\}
\]
are as follows, with a decrement's zero successor written first:
\[
\begin{array}{r|l|l}
\text{site}&P&P_F\\\hline
3550&\operatorname{jump}3553&\operatorname{jump}3553\\
3551&\operatorname{dec}r_1[3553,3552]
    &\operatorname{dec}r_1[3553,3552]\\
3552&\operatorname{jump}3551&\operatorname{jump}3551\\
3553&\operatorname{jump}3637&\operatorname{dec}r_6[3637,3554].
\end{array}
\]
Inspection of the complete resolved instruction graph gives exactly two
edges entering $Z$ from outside:
\[
\begin{array}{r|l|c}
3545&\operatorname{dec}r_6[3551,3546]&3545\longrightarrow3551\\
3548&\operatorname{dec}r_6[3550,3549]&3548\longrightarrow3550.
\end{array}
\]
Both are zero branches of a decrement of $r_6$. This list is an exhaustive
finite graph check, not an observation about a sampled run. The two
programs have exactly the same outside-to-inside edges.
The standalone checker \texttt{verification/four\_site\_cut.py} verifies
the complete instruction lists, their sole difference, and this complete
boundary without importing either source compiler or symbolic simulator.

\begin{lemma}[Four-site zero invariant]
For both programs the assertion
\[
                   \operatorname{pc}\in Z\ \Longrightarrow\ X=0
                                                        \tag{F1}
\]
is preserved by every source step. It holds initially at site 25 for
arbitrary natural values of all nine registers.
\end{lemma}
\begin{proof}
Site 25 is outside $Z$, so the initial assertion has no register
precondition. For preservation, suppose the next program counter lies
in $Z$. If the current counter is outside $Z$, one of the two displayed
zero branches was taken, and therefore $X=0$. If the current counter is
inside $Z$, the induction hypothesis gives $X=0$. Sites 3550, 3551 and
3552 preserve $X$; the old instruction at 3553 preserves it, and the new
instruction there follows its zero branch and preserves it. Thus (F1)
holds after the step. A next counter outside $Z$ satisfies (F1)
immediately.
\end{proof}

\begin{theorem}[Exact preservation of source executions]
For every nine-tuple of natural initial register values, the executions
of $P$ and $P_F$ from source site 25 have identical configurations at
every source time, including identical halting outcomes and halting times.
Here a halting configuration may be made absorbing, or equivalently the
comparison may stop at the first halt. PA proves this statement for the
arithmetized programs.
\end{theorem}
\begin{proof}
Induct on source time. Away from site 3553 the two instructions are
identical. At site 3553, (F1) gives $r_6=0$, and both instructions leave
the entire register vector unchanged and go to site 3637. This proves
agreement of the next complete configurations. Neither program changes
the halt instruction.

The instruction lists, their unique difference, the two boundary edges,
and the four cases in $Z$ are finite data. The proof is induction on a
primitive-recursively coded source run with the formula (F1) as its
invariant. PA has that induction instance and verifies the finite
instruction lookups. In particular, the argument is an ordinary proof
of an internal PA theorem; it does not assume that PA is consistent.
\end{proof}

The source result is stronger than preservation only from the intended
seed. In the machine application the seed is
\[
                   (25;2,0,0,0,0,0,0,0,0).
\]
The invariant is also considerably smaller than a full nine-register
positivity analysis: two entering edges and four source sites suffice.

\subsection{Binding the source proof to the binary table}

The source theorem would not by itself justify editing a binary table.
Three additional bindings are used.

First, the two source resolvers agree on the entire source graph. They
identify the old last emitter token as the jump at site 3553 and the new
last emitter token as the displayed decrement. All other 3,645 source
instructions and all their successor indices agree. The emitter's
formula-building instructions, its guarded gate, and its cleanup are
unchanged.

Second, the literal table is verified on every primitive source outcome.
The new decrement's zero \emph{and positive} outcomes are included, even
though (F1) excludes its positive outcome from the intended source run.
The contracts quantify over arbitrary natural register values, the
variable source gap and an arbitrary finite nonempty spare list. They
produce the complete next source frame and a positive number of physical
steps. Every intermediate control of a nonhalting instruction contract
is nonhalting. The register-list loops used in these contracts are
verified from the literal rows; the affine symbolic presentations encode
arbitrary run lengths, rather than samples of small register values.

Third, exact finite simulation of the delivered table from the all-zero
tape reaches the specified clean frame after 57,768 physical steps,
including the head, all nonzero tape cells, the source address and the
complete physical register/spare file. The source site is 25 and the
logical register vector is the displayed seed. The clean-frame statement
includes the zero exterior of the finite tape support.

\begin{lemma}[Positive finite simulation and converse halting]
PA proves that the delivered binary table halts from blank if and only
if $P_F$ halts from its clean seed.
\end{lemma}
\begin{proof}
Concatenate the positive physical instruction segments after the fixed
startup. A finite source halt supplies a finite literal halt.

For the converse, suppose a literal halting trace has length $T$. If
$T<57{,}768$, it contradicts the checked nonhalting startup by
determinism. Otherwise the startup precedes that halt. Suppose the source
has no halt during its first $T+1$ instructions after the seed. The source
machine's increment, decrement and jump operations are total. Their
physical contracts can therefore be concatenated for all $T+1$ source
instructions. Their lengths are positive, so the resulting physical run
extends beyond time $T$. Its intermediate states are nonhalting. This
contradicts determinism of the given literal halting trace. Thus the
source halts within a bounded finite source trace.

This proof is arithmetical. A source register grows by at most one at each
primitive source instruction; finite source iteration and its encoding
are primitive recursive. The physical contract proofs use terminating
unary sweeps, fixed finite row calculations and induction on the relevant
natural lengths. Concatenation and trace determinism are proved by
induction on coded finite sequences. These are PA proofs, including for
nonstandard trace lengths in a model of PA. No reflection principle or
standard-model soundness argument is used.
\end{proof}

This gives
\[
 \mathrm{PA}\vdash
 H_{\mathrm{PA194F}}\ \longleftrightarrow\ H_{P_F}
                  \ \longleftrightarrow\ H_P.             \tag{F2}
\]
The two binary tables can consume different physical times and spare
lists. Equality of their source configurations is the relevant invariant.

\begin{theorem}[PA194F consistency calibration]
Let $H_{\mathrm{PA194F}}$ be the arithmetic sentence asserting a finite
halting trace of the delivered table from its specified blank start.
For a fixed ordinary primitive-recursive PA proof predicate,
\[
 \mathrm{PA}\vdash H_{\mathrm{PA194F}}\leftrightarrow
    \exists p\,\operatorname{Prf}_{\mathrm{PA}}
                      (p,\ulcorner0=1\urcorner),
\]
and therefore
\[
 \mathrm{PA}\vdash
       \neg H_{\mathrm{PA194F}}\leftrightarrow\operatorname{Con}(\mathrm{PA}).
\]
\end{theorem}
\begin{proof}
Compose (F2) with the proof-producing equivalences in (F3), and negate
the PA-refutation predicate. All transformations and the inductions
establishing their correctness are available in PA. The biconditional
itself does not assume PA consistency; actual standard-model nonhalting
follows under that assumption.
\end{proof}

\paragraph{Precise proof status.}
This is an ordinary mathematical proof of a PA-provability statement,
with the unchanged numerical calibration explicitly attributed to and
retained from the supplied research. The new four-site invariant and
source equivalence are proved above; the literal-table bindings are
supported by complete finite and universal symbolic checks. The latter
use the mathematical soundness of the inherited symbolic rules and
their register-list induction lemmas, as well as the recorded Python
executions; a successful execution is not itself a kernel proof. No new
end-to-end Lean or Rocq theorem is claimed. In particular, neither the
upstream PA196 one-direction formal theorem nor the conditional iPRA
endpoints are relabeled as a kernel-checked proof of this biconditional.
